%% kapa: exam=1 round=33 year=2022 period=1 form=A subject=real_estate subject_ko=부동산학원론 seq=1 number=81
%% source: Q-Net 감정평가사 기출문제 (cst00309) artlSeq=5209997 file=raw/1008/5209997/2022년도 제33회 감정평가사 1차시험 1교시(민법,경제학,부동산학).pdf
%% conversion: kordoc 4.24.51 -> tools/gen1.py (auto); review=proofread-sonnet; flags=none
\begin{question}{81}
토지에 관한 설명으로 옳지 않은 것은?
\begin{choices}
  \item 공간으로서 토지는 지표, 지하, 공중을 포괄하는 3차원 공간을 의미한다.
  \item 자연으로서 토지는 인간의 노력에 의해 그 특성을 바꿀 수 없다.
  \item 소비재로서 토지는 그 가치가 시장가치와 괴리되는 경우가 있다.
  \item 생산요소로서 토지는 그 가치가 토지의 생산성에 영향을 받는다.
  \item 재산으로서 토지는 사용·수익·처분의 대상이 된다.
\end{choices}
\end{question}
\answer{2}
